% !Mode:: "TeX:UTF-8"
% !TEX root = ../../main.tex
% ============================================================
% TOPIC 15
% ============================================================
\mashgulot{15}{Expansion of a vector of a linear space in a basis. Subspaces of a linear space}

\begin{tarif}
A set $V$ on whose elements addition and multiplication by real numbers are
defined is called a \textbf{linear (vector) space} if these operations satisfy
the following eight axioms ($x,y,z\in V$, $\alpha,\beta\in\R$):
\begin{enumerate}[label=\arabic*), itemsep=0pt]
\item $x+y=y+x$; \quad 2) $(x+y)+z=x+(y+z)$;
\setcounter{enumi}{2}
\item $\exists\,\theta\in V$: $x+\theta=x$; \quad 4) $\forall x\ \exists(-x)$: $x+(-x)=\theta$;
\setcounter{enumi}{4}
\item $1\cdot x=x$; \quad 6) $\alpha(\beta x)=(\alpha\beta)x$;
\setcounter{enumi}{6}
\item $(\alpha+\beta)x=\alpha x+\beta x$; \quad 8) $\alpha(x+y)=\alpha x+\alpha y$.
\end{enumerate}
\end{tarif}

Familiar examples are the $n$-dimensional arithmetic space $\R^n$, the space
$P_n$ of polynomials of degree at most $n$, the space $M_{m\times n}$ of
$m\times n$ matrices and the space $C[a;b]$ of functions continuous on $[a;b]$.

\subsection{Linear dependence. Basis and dimension}

\begin{tarif}
Elements $x_1,\dots,x_k$ are called \textbf{linearly dependent} if
$\alpha_1x_1+\dots+\alpha_kx_k=\theta$ for some numbers $\alpha_1,\dots,\alpha_k$
not all equal to zero. Otherwise, that is, if this equality holds only for
$\alpha_1=\dots=\alpha_k=0$, they are called \textbf{linearly independent}.
\end{tarif}

\begin{tarif}
A system $e_1,\dots,e_n$ of elements of $V$ is called a \textbf{basis} if it is
linearly independent and every element of $V$ is a linear combination of
$e_1,\dots,e_n$; the number $n$ is the \textbf{dimension} of the space
($\dim V=n$). The numbers $\xi_1,\dots,\xi_n$ in the expansion
$x=\xi_1e_1+\dots+\xi_ne_n$ are called the \textbf{coordinates} of $x$ in this
basis; they are uniquely determined.
\end{tarif}

\begin{teorema}
Vectors $\vec a_1,\dots,\vec a_n$ of $\R^n$ form a basis if and only if the
determinant made up of their coordinates is nonzero. More generally, the rank
of a system of $k$ vectors equals the dimension of their linear span.
\end{teorema}

\begin{misol}
Show that the vectors $\vec e_1=(1;1;0)$, $\vec e_2=(0;1;1)$, $\vec e_3=(1;0;1)$
form a basis of $\R^3$ and expand $\vec x=(2;3;1)$ in this basis.

\textbf{Solution.} Since
$\begin{vmatrix}1&0&1\\1&1&0\\0&1&1\end{vmatrix}=1(1-0)-0+1(1-0)=2\neq0$,
the vectors form a basis. Writing $\vec x=\xi_1\vec e_1+\xi_2\vec e_2+\xi_3\vec e_3$
coordinatewise, we obtain
\[\kf{\begin{cases}\xi_1+\xi_3=2,\\ \xi_1+\xi_2=3,\\ \xi_2+\xi_3=1.\end{cases}}\]
Adding the three equations gives $\xi_1+\xi_2+\xi_3=3$, from which $\xi_3=0$,
$\xi_1=2$ and $\xi_2=1$. Thus $\vec x=2\vec e_1+\vec e_2$.
\end{misol}

\begin{misol}
Examine the system $\vec a=(1;2;3)$, $\vec b=(2;1;0)$, $\vec c=(3;3;3)$.

\textbf{Solution.} Here $\vec c=\vec a+\vec b$, i.e.\ $\vec a+\vec b-\vec c=\vec 0$,
so the system is linearly dependent. Its rank is $2$, and the linear span is
two-dimensional (a plane through the origin).
\end{misol}

\subsection{Subspaces}

\begin{tarif}
A nonempty subset $L$ of a space $V$ is called a \textbf{subspace} if
$x+y\in L$ and $\alpha x\in L$ for all $x,y\in L$ and $\alpha\in\R$.
The set of all linear combinations of elements $x_1,\dots,x_k$ is called their
\textbf{linear span} and is denoted by $L(x_1,\dots,x_k)$.
\end{tarif}

\begin{misol}
The set $L=\{(x_1;x_2;x_3)\colon x_1+x_2+x_3=0\}$ is a subspace of $\R^3$:
sums and scalar multiples of solutions of a homogeneous equation are again
solutions. Since $x_1=-x_2-x_3$, we can write
$(x_1;x_2;x_3)=x_2(-1;1;0)+x_3(-1;0;1)$; hence the vectors
$(-1;1;0)$, $(-1;0;1)$ form a basis of $L$ and $\dim L=2$.

On the other hand, $M=\{(x_1;x_2)\colon x_1+x_2=1\}$ is not a subspace, because
$(0;0)\notin M$.
\end{misol}

\begin{teorema}
The solutions of a homogeneous system $Ax=0$ form a subspace of $\R^n$ of
dimension $n-\operatorname{rank}A$; any basis of this subspace is called a
\textbf{fundamental system of solutions}.
\end{teorema}

\subsection{Linear spaces and information security}

\textbf{Linear codes.} A binary $(n,k)$ linear code is a $k$-dimensional
subspace of $\mathbb F_2^n$. The rows of the generator matrix $G$ form a basis
of this subspace, and every codeword is a linear combination (mod 2) of the
basis vectors: $c=mG$. The code therefore contains $2^k$ codewords.

\begin{misol}
For $G=\left(\begin{smallmatrix}1&0&1&1\\0&1&1&0\end{smallmatrix}\right)$ the
messages are encoded as $00\to0000$, $10\to1011$, $01\to0110$, $11\to1101$.
These four words form a two-dimensional subspace of $\mathbb F_2^4$.
\end{misol}

\textbf{Parity-check matrix.} A code can equally well be described as the
solution space of the homogeneous system $Hc^T=0$; then
$\dim C=n-\operatorname{rank}H$.

\textbf{Linear cryptanalysis.} The linear part of a cipher is treated as a
linear map of a vector space over $\mathbb F_2$. If the analysis yields $r$
independent linear relations among the $n$ key bits, the admissible keys form
a set of dimension $n-r$ (the solutions of a linear system), and an exhaustive
search shrinks from $2^n$ to $2^{n-r}$ keys.

\subsection{Computing with Python, Maple and Mathcad}

We check the first two examples: the coordinates of $\vec x=(2;3;1)$ in the
basis $\vec e_1,\vec e_2,\vec e_3$, and the rank of the system
$\vec a,\vec b,\vec c$ together with the linear relation between its vectors
(a null vector of the matrix whose columns are $\vec a,\vec b,\vec c$).
\begin{Verbatim}[fontsize=\footnotesize,frame=single,label={\small Python}]
from sympy import Matrix

e1 = Matrix([1, 1, 0])
e2 = Matrix([0, 1, 1])
e3 = Matrix([1, 0, 1])
A = Matrix.hstack(e1, e2, e3)   # basis vectors as columns
print("det =", A.det())
print("coordinates:", A.solve(Matrix([2, 3, 1])).T)

B = Matrix([[1, 2, 3], [2, 1, 0], [3, 3, 3]])  # a, b, c
print("rank =", B.rank())
print("null space of B^T:", B.T.nullspace()[0].T)
\end{Verbatim}
\begin{Verbatim}[fontsize=\footnotesize,frame=single,label={\small Output}]
det = 2
coordinates: Matrix([[2, 1, 0]])
rank = 2
null space of B^T: Matrix([[-1, -1, 1]])
\end{Verbatim}
The last line says that $-\vec a-\vec b+\vec c=\vec 0$. The same results in
Maple:
\begin{Verbatim}[fontsize=\footnotesize,frame=single,label={\small Maple}]
with(LinearAlgebra):
A := Matrix([[1, 0, 1], [1, 1, 0], [0, 1, 1]]):
Determinant(A);                  # 2
LinearSolve(A, <2, 3, 1>);       # <2, 1, 0>
B := Matrix([[1, 2, 3], [2, 1, 0], [3, 3, 3]]):
Rank(B);                         # 2
NullSpace(Transpose(B));         # {<-1, -1, 1>}
\end{Verbatim}
In Mathcad Prime the coordinates are given by \texttt{lsolve} and the rank
by \texttt{rank}:
\begin{Verbatim}[fontsize=\footnotesize,frame=single,label={\small Mathcad}]
     | 1  0  1 |        | 2 |                  | 2 |
A := | 1  1  0 |   x := | 3 |    lsolve(A, x) = | 1 |
     | 0  1  1 |        | 1 |                  | 0 |
|A| = 2
     | 1  2  3 |
B := | 2  1  0 |        rank(B) = 2
     | 3  3  3 |
\end{Verbatim}

\subsection{Exercises on linear spaces and bases}

\emph{1. Expand the vector $\vec x$ in the basis $\vec e_1,\vec e_2,\vec e_3$.}
\bigskip
\begin{enumerate}[label=\arabic*., leftmargin=*]
\item $\vec e_1=(2;\,2;\,-2)$, $\vec e_2=(-1;\,2;\,-1)$, $\vec e_3=(-1;\,0;\,0)$; \ $\vec x=(1;\,-4;\,2)$
\item $\vec e_1=(-1;\,0;\,2)$, $\vec e_2=(1;\,0;\,0)$, $\vec e_3=(2;\,-1;\,1)$; \ $\vec x=(-3;\,-1;\,3)$
\item $\vec e_1=(-1;\,0;\,2)$, $\vec e_2=(0;\,0;\,2)$, $\vec e_3=(-2;\,1;\,2)$; \ $\vec x=(-7;\,2;\,4)$
\item $\vec e_1=(-1;\,-1;\,0)$, $\vec e_2=(-1;\,2;\,1)$, $\vec e_3=(0;\,-2;\,-1)$; \ $\vec x=(0;\,-8;\,-3)$
\item $\vec e_1=(1;\,-1;\,2)$, $\vec e_2=(2;\,-1;\,2)$, $\vec e_3=(-2;\,-1;\,1)$; \ $\vec x=(-1;\,0;\,-1)$
\item $\vec e_1=(0;\,0;\,1)$, $\vec e_2=(-1;\,0;\,2)$, $\vec e_3=(2;\,-1;\,-1)$; \ $\vec x=(5;\,-4;\,0)$
\item $\vec e_1=(-2;\,1;\,2)$, $\vec e_2=(-1;\,1;\,0)$, $\vec e_3=(-1;\,-1;\,1)$; \ $\vec x=(-11;\,2;\,9)$
\item $\vec e_1=(-1;\,-1;\,-1)$, $\vec e_2=(1;\,2;\,2)$, $\vec e_3=(-1;\,-2;\,0)$; \ $\vec x=(0;\,0;\,-2)$
\item $\vec e_1=(2;\,0;\,2)$, $\vec e_2=(2;\,0;\,1)$, $\vec e_3=(0;\,1;\,-2)$; \ $\vec x=(4;\,2;\,1)$
\item $\vec e_1=(2;\,-1;\,-2)$, $\vec e_2=(0;\,-1;\,0)$, $\vec e_3=(-1;\,-2;\,0)$; \ $\vec x=(-2;\,7;\,4)$
\item $\vec e_1=(2;\,1;\,0)$, $\vec e_2=(-2;\,-2;\,2)$, $\vec e_3=(0;\,1;\,-1)$; \ $\vec x=(-6;\,-11;\,12)$
\item $\vec e_1=(0;\,2;\,0)$, $\vec e_2=(1;\,1;\,2)$, $\vec e_3=(0;\,1;\,-1)$; \ $\vec x=(3;\,3;\,6)$
\item $\vec e_1=(2;\,2;\,-1)$, $\vec e_2=(0;\,-1;\,0)$, $\vec e_3=(2;\,2;\,-2)$; \ $\vec x=(-2;\,0;\,4)$
\item $\vec e_1=(-1;\,-1;\,-1)$, $\vec e_2=(1;\,-1;\,0)$, $\vec e_3=(0;\,1;\,0)$; \ $\vec x=(3;\,3;\,3)$
\item $\vec e_1=(1;\,0;\,2)$, $\vec e_2=(1;\,2;\,1)$, $\vec e_3=(1;\,-1;\,2)$; \ $\vec x=(-9;\,0;\,-16)$
\item $\vec e_1=(1;\,-1;\,-2)$, $\vec e_2=(2;\,1;\,-1)$, $\vec e_3=(0;\,-2;\,-1)$; \ $\vec x=(4;\,4;\,-1)$
\item $\vec e_1=(2;\,0;\,-1)$, $\vec e_2=(-1;\,2;\,-1)$, $\vec e_3=(2;\,-2;\,1)$; \ $\vec x=(7;\,-14;\,9)$
\item $\vec e_1=(1;\,1;\,0)$, $\vec e_2=(-1;\,0;\,2)$, $\vec e_3=(-1;\,-1;\,-2)$; \ $\vec x=(0;\,4;\,12)$
\item $\vec e_1=(-1;\,0;\,0)$, $\vec e_2=(-1;\,0;\,2)$, $\vec e_3=(-1;\,1;\,2)$; \ $\vec x=(-2;\,3;\,10)$
\item $\vec e_1=(-1;\,2;\,-1)$, $\vec e_2=(-1;\,-1;\,0)$, $\vec e_3=(1;\,-1;\,0)$; \ $\vec x=(2;\,0;\,-2)$
\item $\vec e_1=(-1;\,2;\,0)$, $\vec e_2=(-1;\,-1;\,-1)$, $\vec e_3=(1;\,2;\,2)$; \ $\vec x=(-9;\,-2;\,-8)$
\item $\vec e_1=(-1;\,2;\,-2)$, $\vec e_2=(1;\,2;\,-1)$, $\vec e_3=(1;\,1;\,-1)$; \ $\vec x=(4;\,16;\,-13)$
\item $\vec e_1=(-2;\,2;\,-1)$, $\vec e_2=(1;\,1;\,1)$, $\vec e_3=(1;\,2;\,2)$; \ $\vec x=(-4;\,1;\,-5)$
\item $\vec e_1=(1;\,2;\,0)$, $\vec e_2=(0;\,1;\,0)$, $\vec e_3=(-2;\,-2;\,1)$; \ $\vec x=(-6;\,-7;\,4)$
\item $\vec e_1=(-2;\,-2;\,0)$, $\vec e_2=(1;\,-1;\,-1)$, $\vec e_3=(0;\,-1;\,0)$; \ $\vec x=(7;\,5;\,-1)$
\item $\vec e_1=(-1;\,-1;\,2)$, $\vec e_2=(-2;\,2;\,2)$, $\vec e_3=(0;\,1;\,0)$; \ $\vec x=(-3;\,13;\,0)$
\item $\vec e_1=(1;\,2;\,-1)$, $\vec e_2=(1;\,-2;\,0)$, $\vec e_3=(2;\,-2;\,-1)$; \ $\vec x=(6;\,-4;\,-4)$
\item $\vec e_1=(-1;\,-2;\,-1)$, $\vec e_2=(0;\,-1;\,-1)$, $\vec e_3=(-1;\,0;\,-1)$; \ $\vec x=(4;\,3;\,5)$
\item $\vec e_1=(1;\,-2;\,-2)$, $\vec e_2=(1;\,0;\,0)$, $\vec e_3=(1;\,-1;\,0)$; \ $\vec x=(4;\,1;\,2)$
\item $\vec e_1=(-2;\,1;\,1)$, $\vec e_2=(-1;\,1;\,0)$, $\vec e_3=(0;\,-1;\,2)$; \ $\vec x=(4;\,-5;\,2)$
\end{enumerate}

\emph{2. Determine whether the system of vectors is linearly independent and find the dimension of its linear span.}
\bigskip
\begin{enumerate}[label=\arabic*., leftmargin=*]
\item $\vec a=(-3;\,-2;\,-2)$, $\vec b=(1;\,1;\,-1)$, $\vec c=(7;\,5;\,3)$
\item $\vec a=(0;\,3;\,-3)$, $\vec b=(-3;\,0;\,-3)$, $\vec c=(-3;\,3;\,-3)$
\item $\vec a=(-1;\,-2;\,-3)$, $\vec b=(-1;\,-3;\,-1)$, $\vec c=(0;\,2;\,-4)$
\item $\vec a=(1;\,0;\,2)$, $\vec b=(1;\,0;\,3)$, $\vec c=(1;\,2;\,-3)$
\item $\vec a=(-2;\,-1;\,-3)$, $\vec b=(2;\,2;\,2)$, $\vec c=(-4;\,-3;\,-5)$
\item $\vec a=(2;\,-1;\,-1)$, $\vec b=(0;\,-1;\,3)$, $\vec c=(2;\,2;\,0)$
\item $\vec a=(1;\,0;\,3)$, $\vec b=(2;\,0;\,2)$, $\vec c=(0;\,0;\,-4)$
\item $\vec a=(1;\,-2;\,3)$, $\vec b=(1;\,-3;\,-3)$, $\vec c=(-2;\,0;\,3)$
\item $\vec a=(1;\,2;\,3)$, $\vec b=(-3;\,3;\,-1)$, $\vec c=(-4;\,10;\,4)$
\item $\vec a=(-2;\,-1;\,0)$, $\vec b=(1;\,0;\,1)$, $\vec c=(2;\,3;\,-1)$
\item $\vec a=(1;\,-2;\,0)$, $\vec b=(-1;\,2;\,0)$, $\vec c=(0;\,0;\,0)$
\item $\vec a=(1;\,3;\,2)$, $\vec b=(1;\,1;\,3)$, $\vec c=(1;\,1;\,-2)$
\item $\vec a=(0;\,-1;\,-3)$, $\vec b=(2;\,2;\,-2)$, $\vec c=(4;\,5;\,-1)$
\item $\vec a=(2;\,-2;\,2)$, $\vec b=(2;\,-1;\,-3)$, $\vec c=(0;\,2;\,3)$
\item $\vec a=(1;\,-2;\,2)$, $\vec b=(-1;\,-1;\,2)$, $\vec c=(2;\,-4;\,4)$
\item $\vec a=(3;\,1;\,-2)$, $\vec b=(3;\,-1;\,1)$, $\vec c=(-3;\,1;\,-2)$
\item $\vec a=(1;\,-1;\,2)$, $\vec b=(-1;\,2;\,2)$, $\vec c=(-1;\,1;\,-2)$
\item $\vec a=(1;\,-3;\,1)$, $\vec b=(0;\,1;\,1)$, $\vec c=(-2;\,3;\,1)$
\item $\vec a=(1;\,2;\,-1)$, $\vec b=(-3;\,0;\,1)$, $\vec c=(-1;\,-2;\,1)$
\item $\vec a=(-2;\,0;\,-2)$, $\vec b=(-1;\,0;\,-2)$, $\vec c=(-1;\,-2;\,2)$
\item $\vec a=(-3;\,3;\,-3)$, $\vec b=(-3;\,2;\,1)$, $\vec c=(-6;\,6;\,-6)$
\item $\vec a=(-2;\,0;\,1)$, $\vec b=(-1;\,-2;\,-1)$, $\vec c=(-1;\,2;\,1)$
\item $\vec a=(-3;\,3;\,-3)$, $\vec b=(0;\,-3;\,3)$, $\vec c=(-6;\,9;\,-9)$
\item $\vec a=(2;\,-1;\,-2)$, $\vec b=(2;\,-2;\,-1)$, $\vec c=(3;\,-3;\,0)$
\item $\vec a=(-1;\,0;\,-2)$, $\vec b=(3;\,0;\,-3)$, $\vec c=(-1;\,0;\,7)$
\item $\vec a=(-3;\,-2;\,-3)$, $\vec b=(-3;\,0;\,2)$, $\vec c=(0;\,3;\,2)$
\item $\vec a=(1;\,0;\,3)$, $\vec b=(-1;\,-3;\,1)$, $\vec c=(1;\,0;\,3)$
\item $\vec a=(0;\,-3;\,3)$, $\vec b=(0;\,-3;\,3)$, $\vec c=(0;\,-1;\,1)$
\item $\vec a=(-1;\,-1;\,-2)$, $\vec b=(3;\,1;\,-1)$, $\vec c=(-2;\,-2;\,-4)$
\item $\vec a=(-3;\,-2;\,-3)$, $\vec b=(-2;\,-1;\,3)$, $\vec c=(3;\,-1;\,1)$
\end{enumerate}

\emph{3. Further problems (10 problems).}
\bigskip
\begin{enumerate}[label=\arabic*., leftmargin=*]
\item \textbf{(A subspace.)} Show that the set $L=\{(x_1;x_2;x_3)\colon x_1+x_2+x_3=0\}$ is a subspace of $\mathbb R^3$ and find a basis of $L$ and its dimension.
\item \textbf{(Not a subspace.)} Is the set $M=\{(x_1;x_2)\colon x_1+x_2=1\}$ a subspace of $\mathbb R^2$?
\item \textbf{(Membership in a span.)} Does the vector $\vec x=(4;1;5)$ belong to the linear span of $\vec a=(1;-1;2)$ and $\vec b=(2;3;1)$?
\item \textbf{(Membership in a span.)} Does the vector $\vec x=(1;2;4)$ belong to the linear span of $\vec a=(1;0;1)$ and $\vec b=(0;1;1)$?
\item \textbf{(A space of polynomials.)} In $P_2$, expand $p(x)=x^2+3x+2$ in the basis $1,\ x-1,\ (x-1)^2$.
\item \textbf{(A space of matrices.)} What is the dimension of the space of $2\times2$ matrices? Write down its standard basis.
\item \textbf{(Change of basis.)} Find the coordinates of $\vec x=(5;1)$ in the basis $\vec e_1=(1;1)$, $\vec e_2=(1;-1)$.
\item \textbf{(A binary space.)} Are the vectors $\vec a=(1;1;0)$, $\vec b=(0;1;1)$, $\vec c=(1;0;1)$ of $\mathbb{F}_2^3$ linearly independent (over $\mathbb F_2$, i.e.\ mod 2)?
\item \textbf{(A homogeneous system.)} Find a basis and the dimension of the solution space of the system $\begin{cases}x_1+2x_2-x_3=0\\2x_1+4x_2-2x_3=0\end{cases}$.
\item \textbf{(A code subspace.)} How many codewords does the binary code with generator matrix $G=\left(\begin{smallmatrix}1&0&1&1\\0&1&1&0\end{smallmatrix}\right)$ have? List them.
\end{enumerate}

